% Bewerbungs-Template: Preamble
% Wird von templates/bewerbung.tex eingebunden.

\documentclass[10.5pt,a4paper]{scrartcl}

% Keine automatische Sektionsnummerierung -- wir wollen kein "1.1 Profile"
\setcounter{secnumdepth}{-2}

% --- Sprache (de oder en, wird aus meta.tex via \metaspracheXX gesetzt) ---
\usepackage{iftex}
\usepackage[ngerman,english]{babel}

% --- Font via fontspec (lualatex / xelatex) ---
% Font family + weight names come from meta.tex (generated by yaml_to_tex.py
% from user.yaml). Default is "TeX Gyre Heros" (bundled with scheme-full).
% Override in user.yaml:
%   font_family: "Fira Sans"
%   font_features: "Numbers=OldStyle, Ligatures={Common,TeX}, ..."
%
% Ligatures=TeX enables the classic TeX-glyph ligatures:
%   --   -> en-dash (-)          ---  -> em-dash (--)
%   ``''  -> typographic quotes
%   <<>>  -> French guillemets
% Token-level the source stays unchanged (\IfSubStr{...}{--} still works in
% \cvsubsection), only the rendered glyph is substituted.
\usepackage{fontspec}
\setmainfont{\metafontregular}[\metafontfeatures]

\newfontfamily\fontsemibold{\metafontsemiboldname}[\metafontfeatures]
\newfontfamily\fontmedium{\metafontmediumname}[\metafontfeatures]
\newfontfamily\fontregular{\metafontregular}[\metafontfeatures]
\newfontfamily\fontitalic{\metafontitalicname}[\metafontfeatures]


% --- REM ---
\newlength{\rem}
\setlength{\rem}{0.4cm}

% --- Packages ---
\usepackage{fontawesome5}
\usepackage{graphicx}
\usepackage{tikz}        % fuer das Foto mit hartem Schatten (document.tex)
\usepackage{ragged2e}    % \RaggedRight im Anschreiben
\usepackage{microtype}
% Hanging punctuation: Satzzeichen am Zeilenanfang/-ende ragen leicht in den
% Margin, damit der optische Rand sauber bleibt (Knuth: "protrusion"). Default
% in microtype lualatex aktiv -- explizit nur zur Sicherheit.
\microtypesetup{protrusion=true,activate={true,nocompatibility}}
\usepackage[normalem]{ulem}
\usepackage[all]{nowidow}
\usepackage{xstring}     % \IfSubStr / \StrBefore / \StrBehind -- fuer \cvsubsection, \cvdlitem

% --- Seitenlayout: angelehnt an das ODT (oben 0.85cm, unten 1.5cm, links 2.5cm, rechts 2cm) ---
\usepackage[
  a4paper,
  top    = {\dimexpr\rem*4\relax},
  bottom = {\dimexpr\rem*5\relax},
  left   = {\dimexpr\rem*5\relax},
  right  = {\dimexpr\rem*5\relax},
  headheight = {\dimexpr\rem*3\relax},
  headsep    = \rem,
  footskip   = {\dimexpr\rem*2\relax},
]{geometry}

% --- Farben (zurueckhaltend) ---
\usepackage{xcolor}
\definecolor{text}{HTML}{282A36}
\definecolor{muted}{HTML}{44475A}
\definecolor{purple}{HTML}{BD93F9}


% --- Listen, Bilder, Layout-Helfer ---
\usepackage{enumitem}
\setlist[itemize]{
  leftmargin = \dimexpr 0.6\rem\relax,
  itemsep    = \dimexpr 0.2\rem\relax,
  parsep     = 0pt,
  topsep     = 0em,
label      = \textcolor{purple}{\faAngleRight[regular]}\hspace{\dimexpr -0.3\rem\relax},
}

% Asset-Pfad: erst absoluter Container-Mount, dann relativer Fallback fuer
% lokale Builds ausserhalb des Containers. LaTeX probiert in Reihenfolge.
\graphicspath{{/job/assets/}{../../assets/}}

% --- Kopf-/Fusszeile: Seite X von Y ---
\usepackage{lastpage}
\usepackage{scrlayer-scrpage}
\pagestyle{scrheadings}
% Footer linksbuendig (\lofoot* statt \cfoot*). Anschreiben nutzt
% \thispagestyle{plain.scrheadings} -> kein Footer. CV setzt page-Counter
% bei 1 neu und nutzt \pagestyle{scrheadings}.
\clearpairofpagestyles
\lofoot*{\color{muted}\fontsizeabsender\metaseite~\thepage~\metavon~\pageref{LastPage}}
\setkomafont{pagefoot}{\fontregular}

% --- Hyperlinks (dezent, ohne Box) ---
\usepackage[
  colorlinks = true,
  urlcolor   = text,
  linkcolor  = text,
  citecolor  = text,
  pdfborder  = {0 0 0},
]{hyperref}


% --- Unterstreichungs-Einstellungen ---
\renewcommand{\ULdepth}{\dimexpr 0.35\rem\relax}        % Abstand der Linie zur Schrift
\renewcommand{\ULthickness}{\dimexpr 0.2\rem\relax}    % Strichstaerke

% In \href einbauen
\let\oldhref\href
\renewcommand{\href}[2]{\oldhref{#1}{#2}}

\newcommand{\bodylink}[2]{\oldhref{#1}{\textcolor{purple}{\uline{\textcolor{text}{#2}}}}}

% Icon-Link: kein \uline, nur purple einfaerben (fuer FontAwesome-Icons)
\newcommand{\iconhref}[2]{\oldhref{#1}{\textcolor{purple}{#2}}}

% --- Absatz-Stil: kein Einzug, Abstand zwischen Absaetzen ---
\setlength{\parindent}{0pt}
\setlength{\parskip}{\glueexpr 1.0\rem plus 0.4\rem\relax}

% Seiten unten *nicht* per parskip-Stretch ausgleichen -- sonst verteilt TeX
% den ungenutzten Platz auf die Absatz-Abstaende, und kompaktere Aenderungen
% in einer Sektion (z.B. Skills) verschieben sichtbar Abstaende in *anderen*
% Sektionen. \raggedbottom laesst die Seite unten ruhig kuerzer enden.
% \raggedbottom

% --- Schriftgrößen ---
\newcommand{\fontsizebody}{\fontsize{\dimexpr 0.9\rem\relax}{\dimexpr 1.2\rem\relax}\selectfont}
\newcommand{\fontsizebetreff}{\fontsize{\dimexpr 1.4\rem\relax}{\rem}\selectfont}
\newcommand{\fontsizeheader}{\fontsize{\dimexpr 1.2\rem\relax}{\rem}\selectfont}
\newcommand{\fontsizesubheader}{\fontsize{\dimexpr 1\rem\relax}{\rem}\selectfont}
\newcommand{\fontsizeabsender}{\fontsize{\dimexpr 0.9\rem\relax}{\dimexpr 0.7\rem\relax}\selectfont}

% --- Section-Schriftgrade & Abstaende dezenter ---
\addtokomafont{section}{\fontsemibold\fontsizebetreff\color{text}}
\addtokomafont{subsection}{\fontsemibold\fontsizeheader\color{text}\hspace{\dimexpr -0.6\rem\relax}}
\addtokomafont{subsubsection}{\fontsemibold\fontsizesubheader\color{text}}
% Hinweis: scrartcl kennt nur bis \subsubsection. Tieferes Level (####) wird
% ueber den Markdown-Renderer headingFour -> \cvsubsection abgefangen, der
% intern \subsubsection ruft -- daher kein \addtokomafont{subsubsubsection} noetig.
\RedeclareSectionCommand[
  beforeskip = 0.0cm,
  afterskip  = 0.0cm,
  runin      = false,
]{section}
\RedeclareSectionCommand[
  beforeskip = 0.0cm,
  afterskip  = \dimexpr -0.4\rem\relax,
  runin      = false,
]{subsection}
\RedeclareSectionCommand[
  beforeskip = 0.0cm,
  afterskip  = \dimexpr -0.8\rem\relax,
  runin      = false,
]{subsubsection}

% --- Helper-Makros ---

% Sprachabhaengige Labels (gesteuert ueber \metaspracheen aus meta.tex)
\newcommand{\cvtitle}{\ifdefined\metaspracheen Curriculum Vitae\else Lebenslauf\fi}
\newcommand{\anlagenlabel}{\ifdefined\metaspracheen Enclosures\else Anlagen\fi}

% Save-Box fuer Profil-Layout (Bildhoehe = Profil-Boxhoehe, zwei-Spalten)
\newsavebox{\profilbox}
\newsavebox{\imgbox}     % fuer Foto + Schatten-Overlay in document.tex

% Spaltenbreite fuer das Foto neben dem Profil (Skip = stretch/shrink moeglich).
\newlength{\bildbreite}
\setlength{\bildbreite}{\glueexpr 8\rem plus 2\rem\relax}

% --- Flex-Layout: CSS-flexbox-flow ohne Spalten -----------------------------
% Jedes \flexitem wird in seiner natuerlichen Breite gesetzt (\sbox), gefolgt
% von einem festen \flxgap. Wenn das naechste Item nicht mehr in die aktuelle
% Zeile passt, kommt ein \par (neue Zeile). Items mit natuerlicher Breite
% groesser als \linewidth werden in einer parbox umbrochen.
%
% Aufruf:
%   \flexstart
%   \flexitem{kurz}\flexitem{mittel langer text}\flexitem{wirklich sehr lang}
%   \flexend
\newsavebox{\flxbox}
\newdimen\flxitemwd
\newdimen\flxrowused
\newdimen\flxgap
\newdimen\flxrowsep
\newdimen\flxcatabove
\newdimen\flxcatbelow
\flxgap     = \dimexpr 0.6\rem\relax   % horizontaler Abstand zwischen Items
\flxrowsep  = \dimexpr 0.1\rem\relax   % vertikaler Abstand zwischen Reihen
\flxcatabove= \dimexpr -0.4\rem\relax   % Abstand ueber Kategorie-Heading
\flxcatbelow= \dimexpr -0.8\rem\relax   % Abstand unter Kategorie-Heading (vor items)

\newcommand{\flxbullet}[1]{%
  \textcolor{purple}{\faAngleRight[regular]}\hspace{\dimexpr 0.2\rem\relax}#1%
}

% \flexstart oeffnet eine Gruppe und ueberschreibt \parskip lokal mit
% \flxrowsep, sodass der vertikale Abstand zwischen Reihen kontrolliert
% wird (sonst wuerde der globale \parskip = 1.2rem dazwischen greifen).
\newcommand{\flexstart}{%
  \par\noindent
  \begingroup
  \setlength{\parskip}{\flxrowsep}%
  \flxrowused=0pt
  \noindent
}

\newcommand{\flexitem}[1]{%
  \sbox\flxbox{\flxbullet{#1}}%
  \flxitemwd=\wd\flxbox
  % Item breiter als \linewidth -> in parbox umbrechen, full width
  \ifdim\flxitemwd>\linewidth
    \sbox\flxbox{\parbox[t]{\linewidth}{\flxbullet{#1}}}%
    \flxitemwd=\linewidth
  \fi
  % Row-break wenn das Item nicht mehr in die aktuelle Zeile passt
  \ifdim\dimexpr\flxrowused+\flxitemwd\relax>\linewidth
    \par\noindent
    \flxrowused=0pt
  \fi
  \usebox\flxbox
  \hspace{\flxgap}%
  \advance\flxrowused by \dimexpr\flxitemwd+\flxgap\relax
}

\newcommand{\flexend}{\par\endgroup}

% Skill-Kategorie-Heading (innerhalb skills.tex vor jedem \flexstart-Block).
% Abstaende ueber/unter dem Heading via \flxcatabove / \flxcatbelow.
\newcommand{\flexcategory}[1]{%
  \par\vspace{\flxcatabove}%
  {\fontsemibold\fontsizesubheader\color{text}#1}\par
  \vspace{\flxcatbelow}%
}

% Hilfsmakro: skills.tex einlesen mit \endlinechar=-1, sonst injiziert
% das markdown-Renderer-Umfeld zwischen Zeilen \par-Trenner und die
% \flexitem-Reihen brechen vor der eigentlichen Row-Logik.
\newcommand{\inputskills}{%
  \begingroup\endlinechar=-1\relax
  % Skills in flex-layout. Pick categories per application: keep what matches
% the posting, drop the rest. Items within a category: posting keywords in
% their exact phrasing first.
%
% Format:
%   \flexcategory{Category heading}
%   \flexstart
%     \flexitem{Item}
%     \flexitem{Another item}
%   \flexend

\flexcategory{Category A}
\flexstart
  \flexitem{Item}
  \flexitem{Item}
  \flexitem{Item}
\flexend

\flexcategory{Category B}
\flexstart
  \flexitem{Item}
  \flexitem{Item}
\flexend

\flexcategory{Category C}
\flexstart
  \flexitem{Item}
  \flexitem{Item}
\flexend
%
  \endgroup
}

% CV-Job-Subsubsection: wenn der Heading "Firma -- Datum" enthaelt, links/rechts
% mit \hfill setzen; sonst normaler subsubsection. Wird vom markdown-Renderer
% headingThree aufgerufen (siehe \markdownSetup unten).
\newcommand{\cvsubsection}[1]{%
  \IfSubStr{#1}{--}{%
    \StrBefore[1]{#1}{ --}[\jobfirma]%
    \StrBehind[1]{#1}{-- }[\jobdatum]%
    \subsubsection{\makebox[\linewidth][s]{%
    {\fontmedium\fontsizebody\color{muted} \jobfirma}\hfill{\fontmedium\fontsizebody\color{muted}\jobdatum}%
    }}%
  }{%
    \subsubsection{\fontsemibold\fontsizeheader{#1}}%
  }%
}

% Definition-List-Item (Projects): nur Titel (fett), optional inline-URL
% nach dem Titel als GitHub-Icon. Kein rechtsbuendiger Untertitel mehr --
% der Text nach " -- " wird ignoriert (Backward-Compat alter Markdown-Files).
%   "Titel"                       -> Titel
%   "Titel https://URL"           -> Titel + GitHub-Icon (Link)
%   "Titel https://URL -- ..."    -> wie oben (Rest verworfen)
% Description (": ...") bleibt unveraendert darunter.
\newcommand{\cvdlitem}[1]{%
  \par\vspace{\dimexpr -0.6\rem\relax}%
  \noindent
  \llap{\textcolor{purple}{\faArrowRight[solid]}\hspace{\dimexpr 0.5\rem\relax}}%
  \IfSubStr{#1}{--}{%
    \StrBefore[1]{#1}{ --}[\dlfirst]%
  }{%
    \def\dlfirst{#1}%
  }%
  \IfSubStr{\dlfirst}{https://}{%
    \StrBefore[1]{\dlfirst}{ https://}[\dltitle]%
    \StrBehind[1]{\dlfirst}{https://}[\dlurl]%
    {\fontsemibold\fontsizebody\dltitle}\hspace{\dimexpr 0.3\rem\relax}%
    \iconhref{https://\dlurl}{\faGithub}%
  }{%
    {\fontsemibold\fontsizebody\dlfirst}%
  }%
  \par\vspace{\dimexpr -1.0\rem\relax}%
}

% Absender-Zeile (einzeilig): Name + Adresse links, drei FontAwesome-Icons
% (Telefon, Email, GitHub) rechts via \hfill -- alle drei klickbar.
\newcommand{\absenderzeile}{%
  \fontsizeabsender
  {\fontsemibold\metaname}\ \textcolor{purple}{\faAngleRight[regular]}\ \fontsizeabsender\metaadresse\ \textcolor{purple}{\faAngleRight[regular]}\ \metaplzort
  \par
  \vspace{\dimexpr -0.9\rem\relax}
  \hspace{\dimexpr -1.2\rem\relax}
\textcolor{purple}{\fontsize{\dimexpr 0.7\rem\relax}{\rem}\selectfont\faPhone[solid]}\hspace{\dimexpr 0.3\rem\relax}\href{\metatelhref}{\textcolor{muted}{\metatel}}
  \hspace{\dimexpr 0.2\rem\relax}
\textcolor{purple}{\fontsize{\dimexpr 0.7\rem\relax}{\rem}\selectfont\faEnvelope[solid]}\hspace{\dimexpr 0.3\rem\relax}\href{mailto:\metaemail}{\textcolor{muted}{\metaemail}}
  \hspace{\dimexpr 0.2\rem\relax}
\textcolor{purple}{\fontsize{\dimexpr 0.9\rem\relax}{\rem}\selectfont\raisebox{\dimexpr -0.1\rem\relax}{\faGithub}}\hspace{\dimexpr 0.3\rem\relax}\href{https://\metagithub}{\textcolor{muted}{\metagithub}}
}

% --- Markdown-Renderer (witiko/markdown CTAN package) ---
% Erlaubt \markdownInput{file.md} und ersetzt die pandoc-Pipeline.
% Vollstaendige Options-Liste: `texdoc markdown`.
\usepackage[
  smartEllipses = true,   % "..." -> echte Ellipsis
  fencedCode    = true,   % ``` code blocks
  pipeTables    = true,   % | tabellen |
  definitionLists = true,
  % hybrid bewusst aus: damit werden &, %, # etc. automatisch escaped.
  % Wer LaTeX-Befehle im Markdown braucht, kann sie in \begin{markdown}-Bloecke
  % als raw-tex einbetten oder das Macro pro Stelle aktivieren.
]{markdown}

\markdownSetup{
  renderers = {
    headingOne        = {\section{#1}},
    headingTwo        = {\subsection{#1}}, % "## ..." -> \subsection (Skills wird ueber document.tex eingebunden)
    headingThree      = {\subsubsection{#1}},                   % alle ### (z.B. Rolle, Projekt-Titel)
    headingFour       = {\cvsubsection{#1}\vspace{\dimexpr 0.1\rem\relax}},                    % alle #### (splittet bei " -- " auf Firma\hfill Datum)
    strongEmphasis    = {\fontmedium{#1}}\fontregular,   % alle **fett**
    emphasis          = {\fontitalic{#1}}\fontregular,                       % alle _kursiv_
    link              = {\bodylink{#2}{#1}},                     % alle [text](url)
    dlBegin           = {\vspace{\dimexpr 0.6\rem\relax}},
    dlEnd             = {},
    dlItem            = {\cvdlitem{#1}},
    dlItemEnd         = {},
    dlDefinitionBegin = {\noindent\fontsizebody},
    dlDefinitionEnd   = {\par},
  },
}
